\section{General Structure}
The root folder contains various directories, specifically:
\begin{itemize}
    \item analysis
    \item archive
    \item Benchmark Framework
    \item results
\end{itemize}

\subsection{analysis}
This folder contains post run analysis material, such as notebooks and derived files used to inspect and summarize the benchmark results. Raw outputs, final reports and slides are stored elsewhere.

\subsection{archive}
This folder contains material provided by the professor, including research papaers, a list of PDDl planning domains, two reference Github repositories and a Quick PDDL guide that we wrote to better understand the main features of the language.  

\subsection{Benchmark Framework structure}
The Benchmark Framework folder is the core of our work. The main goal of this framework was to provide an easy-to-use and general testing platform to be employed for evaluating the performance of LLMs on PDDL planning tasks. 
The following pages describe the folder structure of the framework and explain the role of its main components:

\vspace{0.2cm}

\begin{center}
        

    \begin{forest} 
            for tree={
                font=\ttfamily\footnotesize,
                grow'=0,
                child anchor=west,
                parent anchor=south,
                anchor=west,
                calign=fixed edge angles,
                edge={draw, semithick},
                l sep = 7pt,       % Ridotto lo spazio orizzontale tra i livelli
                s sep = 3pt,       % Ridotto lo spazio verticale tra le cartelle
                inner sep = 2.5pt,  % ridotto spazio attorno al testo
                folder, % stile cartella
            }
            [LLM Benchmark/
                [docs/]
                [evaluators/]
                [Leonardo script/]
                [models/]
                [outputs/]
                [prompts/]
                [protocols/]
                [reporting/]
                [runner/]
                [scripts/]
                [slurm logs/]
                [tasks/]
                [tests/]
                [utils/]
                [advanced planning evaluation.py]
                [clear outputs.py]
                [run benchmark.py]
                [setup.md]
            ]
            \end{forest}
\end{center}

\paragraph{\texttt{docs}}
Contains operational documentation that is too specific for the main README or setup guide. In particular, it includes documentation for preparing Hugging Face model weights before running local on HPC. This is important because GPU benchmark jobs should not download large model weights during execution: models should be prepared first and then loaded from local files.

\paragraph{\texttt{evaluators}}
Provides the shared evaluation layer used by all model backends and protocols. It includes tools for parsing model outputs, validating generated plans with VAL, classifying errors and computing comparable metrics across runs.

\paragraph{\texttt{Leonardo_script}}
Contains SLURM scripts used to launch, configure, and manage benchmark jobs on the Leonardo HPC system. These scripts handle environment activation, model preparation, cache checks and benchmark execution on compute nodes.

\paragraph{\texttt{models}}
Manages model registries and backend adapters. The adapters provide a common interface for different execution backends, such as the NVIDIA API, local Hugging Face models.

\paragraph{\texttt{outputs}}
Contains the generated benchmark artifacts. The outputs are divided into subfolders such as \textit{raw}, \textit{parsed}, \textit{scored}, and \textit{logs}. The \textit{raw} folder stores prompts and model responses, \textit{parsed} stores the extracted and structured plans, and \textit{scored} stores validation results, metrics and benchmark summaries.

\paragraph{\texttt{prompts}}
Contains the text used to build the prompts given to the LLMs. These files define the general benchmark rules, task-specific instructions, examples and feedback templates used during iterative repair.

\paragraph{\texttt{protocols}}
Defines the different benchmark execution strategies. For example, \texttt{direct_plan} asks the model to solve the problem in a single attempt, while \texttt{iterative_repair} allows multiple attempts using validator feedback.

\paragraph{\texttt{reporting}}
Contains post-run analysis helpers used to generate plots and additional reports from the benchmark outputs. These scripts are used after the benchmark has finished, not during the main execution loop. This folder could be useful for future implementations.

\paragraph{\texttt{requirements}}
Contains the Python dependency files used to configure the environments, especially for Leonardo HPC runs. Different requirement files are used for different model setups.

\paragraph{\texttt{runner}}
Contains the orchestration logic used to launch a single benchmark case or an entire benchmark suite. It connects tasks, prompts, protocols, model adapters, parsing, validation, metrics and output writing.

\paragraph{\texttt{scripts}}
Contains utility scripts that support benchmark operation. For example, it includes scripts for preparing Hugging Face models before execution and for computing domain complexity scores.

\paragraph{\texttt{slurm_logs}}
Stores stdout and stderr logs generated by SLURM jobs on Leonardo. These logs are mainly used for debugging HPC execution, environment setup and job-level failures.

\paragraph{\texttt{tasks}}
Contains the PDDL domains and problem instances used by the benchmark. The instances are divided into three difficulty levels: \textit{easy}, \textit{medium} and \textit{hard}.

\paragraph{\texttt{tests}}
Contains unit and integration-style tests used to verify the benchmark pipeline before running real experiments. These tests check components such as parsing, validation, model adapters, suite execution, and output generation.

\paragraph{\texttt{utils}}
Contains bundled VAL binaries and related validation utilities for Linux and Windows. These executables are used to validate PDDL plans during benchmark execution.

\paragraph{\texttt{advanced_planning_evaluation_sp.py}}
Reads benchmark artifacts from the \texttt{outputs} folder, computes additional analysis metrics and writes a model-centric JSON report under the \texttt{results} folder.

\paragraph{\texttt{clear_outputs.py}}
Safely clears generated benchmark outputs while preserving the folder structure and required placeholder files, used to clean folders during benchmark execution on Leonardo HPC.

\paragraph{\texttt{run_benchmark.py}}
Is the main entry point used to run the benchmark suite. It parses command-line arguments, selects models, protocols, tasks, validators, output settings and then launches the benchmark execution pipeline.

The framework is organized in this way to keep all the logic clearly separated. This modular structure makes it easier to modify important parameters, such as the selected models, protocols, task families, difficulty levels, or number of repair iterations and then rerun the benchmark with the updated configuration.

\subsection{results}
This folder contains all the plots and the output archive of all the analysis process; the outputs folder contains only raw datas the results folder contains the analysis done on the data. 